\documentclass[10.5pt]{article}
\usepackage[margin=0.85in]{geometry}
\usepackage{amsmath,amssymb,mathtools}
\usepackage{booktabs,array,tabularx,longtable}
\usepackage[colorlinks=true,urlcolor=blue,linkcolor=blue]{hyperref}
\setlength{\parindent}{0pt}\setlength{\parskip}{3.5pt}
\newcommand{\RPS}{\mathrm{RPS}}
\newcommand{\ra}{\raggedright\arraybackslash}
\begin{document}

\begin{center}
{\Large\bfseries A non-circular evaluation protocol for CE-trained vs correlation-trained judges}\\[4pt]
Dated 11-21-2026. Written before any new data were looked at.
\end{center}

\textbf{By:} ultimate\_brando9 (Brando's AI assistant), Claude Sonnet 5.5, Claude Code headless harness. No model was called and no experiment was run for this note. The only computation is a 40-line numeric illustration of the algebra (Section 2, script \texttt{toy\_circular.py} in the same folder). \textbf{Project:} \texttt{judges-dont-need-psychometrics-or-corr}. \textbf{Builds on:} experiment 18 (maths R1--R7, W1--W7), experiment 20 (limited-data results), experiment 21 (Codex's frozen fresh-family protocol, commit ff2f07a), and Sanmi's feedback (saved in the repo, commit 10897ca).

\section*{Bottom line (11-21-2026)}
\begin{enumerate}
\item \textbf{Scoring a Pearson-trained judge with Pearson correlation is circular twice.} (a) The training objective and the test metric are the same function of the judge's mean ratings, so the Pearson-trained judge is favoured by construction. (b) That function sees only the mean rating per item, up to a shift and a positive rescaling (R1, R2), so it ties or inverts judges that differ a lot in level, scale, spread and shape. In the toy of Section 2, a judge with $\rho=1.000$ has ranked probability score (RPS) 0.189, no better than the uniform forecast (0.192), while a judge with $\rho=0.949$ has RPS 0.108.
\item \textbf{Primary metrics:} held-out ranked probability score (RPS), then cross-entropy (CE) and total variation distance (TVD), each against the human rating distribution, always next to the uniform forecast. \textbf{Metrics neither objective optimises:} top-1 agreement with the human modal rating, calibration of the forecast cumulative probabilities, and calibration slope of human mean on judge mean. Pairwise preference accuracy is added for ranking uses. Pearson, Spearman and Lin's concordance correlation are secondary columns.
\item \textbf{Decision rule:} find the smallest number $N^\ast$ of human ratings at which the CE-trained judge beats both the Pearson-trained judge and the uniform forecast on held-out RPS, with an interval, over held-out task families. If the rating budget $B$ buys at least $N^\ast$ ratings ($B/c\ge N^\ast$, with $c$ the cost of one rating), train with the proper loss; below $N^\ast$ do not fit a flexible head by either loss.
\item \textbf{Both ways stay open.} Where Pearson wins (ranking uses, picking one judge among a few, very small $N$ as in experiment 20) the protocol reports it in the same table.
\end{enumerate}

\section*{1. Limits of each metric (read this first)}
{\small\renewcommand{\arraystretch}{1.15}
\begin{longtable}{@{}>{\ra}p{2.5cm}>{\ra}p{3.7cm}>{\ra}p{3.9cm}>{\ra}p{1.5cm}>{\ra}p{3.3cm}@{}}
\toprule
Metric & Reads & Blind to & Proper for $p^\ast$? & Favoured by construction\\
\midrule
Pearson $\rho$ between judge means and human means & one mean $m_i$ per item, up to a shift and a positive rescale (R1, R2) & level, scale, spread, skew, all shape & no & the Pearson-trained judge (its own objective)\\
Spearman $\rho_s$ & ranks of the means & same, plus the size of gaps & no & none directly; rank version of Pearson\\
Lin's concordance (CCC) & means with level and scale: $\dfrac{2s_{mt}}{s_m^2+s_t^2+(\bar m-\bar t)^2}$ & spread, shape & no & none; penalises level and scale that Pearson ignores\\
\textbf{RPS} (primary) & the whole forecast CDF at all $K-1$ thresholds & little; noisy when each item has few ratings & yes, strictly & the CE-trained judge in the population limit (same minimiser $p^\ast$); not its training loss\\
\textbf{CE} & each category probability & very large penalties on near-zero probabilities, so the probability floor $\varepsilon$ matters (Sanmi) & yes, strictly & the CE-trained judge (its own objective)\\
\textbf{TVD} & $L_1$ distance to the empirical rating histogram & with $k$ ratings per item the expected score is biased (with $k=1$ it is $1-\sum_r p^\ast(r)p(r)$, R6) & no for sparse labels & none\\
Top-1 agreement & the forecast's modal rating & everything except the mode & no & none\\
Calibration (reliability, slope) & whether forecast cumulative probabilities and means match what raters did & ranking quality & no & none\\
Pairwise preference accuracy & the sign of mean differences between item pairs & level, scale, spread (it reads means only, like Pearson) & no & none directly\\
\bottomrule
\end{longtable}}

\section*{2. Why Pearson scoring of Pearson-trained and CE-trained judges is circular}
\textbf{Notation.} Items $i=1,\dots,n$; ratings $r\in\{0,\dots,K-1\}$; the judge forecasts a distribution $p_i$ over ratings; humans have a distribution $p^\ast_i$ (observed through $k_i$ ratings $y_{i1},\dots,y_{ik_i}$). Means $m_i=\sum_r r\,p_i(r)$ (judge) and $t_i=\sum_r r\,p^\ast_i(r)$ (humans), averages $\bar m,\bar t$, and $\tilde t_i=t_i-\bar t$. $F_q(k)=\sum_{r\le k}q(r)$ is the cumulative distribution. Pearson: $\rho(m,t)=\dfrac{\sum_i(m_i-\bar m)(t_i-\bar t)}{\sqrt{\sum_i(m_i-\bar m)^2\sum_i(t_i-\bar t)^2}}$.

\textbf{Step 1 (R1: Pearson reads only the means).} The numerator is $\sum_i(m_i-\bar m)(t_i-\bar t)=\sum_i\tilde t_i m_i-\bar m\sum_i\tilde t_i=\sum_i\tilde t_i m_i$, because $\sum_i\tilde t_i=0$. The denominator, for a fixed target $t$, is a function of $m$ alone. So $\rho=F_t(m_1,\dots,m_n)$: two judges with the same means have the same $\rho$ whatever their $p_i$.

\textbf{Step 2 (R2: shift and positive rescale do not matter).} For $b>0$, $\rho(a+bm,t)=\rho(m,t)$: the shift cancels in the centering, and $b$ multiplies numerator and denominator. So every judge with $m_i=a+bt_i$, $b>0$, has $\rho=1$.

\textbf{Step 3 (the objective is the metric).} The Pearson-trained judge maximises $\rho$ on training items, i.e.\ $F_t$ evaluated on the training means. Scoring it on test items with Pearson evaluates the same functional $F_t$ on new items. Empirical risk minimisation of a metric on training data is expected to win on that metric on test data; this says nothing about whether the output is useful as a rating or a distribution. A fair test needs a metric that is not the training objective of exactly one arm, or one that no arm optimises.

\textbf{Step 4 (the metric cannot see what CE learns).} CE and RPS are proper: their expected value over $Y\sim p^\ast$ is minimised only at $p=p^\ast$ (R6), which fixes the mean and every higher moment. Pearson gives the same value $1$ to $p=p^\ast$ and to any judge with $m=a+bt$, however wrong its level, scale, spread and shape.

\textbf{Numeric illustration} (exact expectations under the human distributions, not an experiment). $K=6$, five items with human means $t=(0.5,1.5,2.5,3.5,4.5)$, human distributions $\text{Binomial}(5,t_i/5)$; expected scores under those distributions, lower is better for RPS, CE, TVD. Judge A compresses the means, $m_i=2.5+0.1(t_i-2.5)$. Judge B has means $t_i+e_i$ with $e=(+0.5,-0.5,+0.5,-0.5,0)$.

{\small\begin{center}\begin{tabular}{@{}lccccc@{}}
\toprule
Forecast & Pearson $\rho$ (higher better) & RPS & CE (nats) & TVD & rank by Pearson / by RPS\\
\midrule
Truth $p^\ast$ & 1.000 & 0.0953 & 1.234 & 0.000 & 1 (tie) / 1\\
A: compressed means, $\rho=1$ & 1.000 & 0.1891 & 1.983 & 0.397 & 1 (tie) / 3\\
B: true means plus errors & 0.949 & 0.1083 & 1.341 & 0.169 & 3 / 2\\
Uniform & undefined (constant means) & 0.1922 & 1.792 & 0.427 & -- / 4\\
\bottomrule
\end{tabular}\end{center}}
Pearson ranks A above B; RPS ranks B far above A, and A is no better than uniform on RPS and worse than uniform on CE. This is the algebra above in numbers, not evidence about real judges.

\textbf{The symmetric caveat.} RPS is not neutral either: the CE-trained judge is consistent for $p^\ast$, which is what RPS measures. That is the substance of the claim (``read the output as a distribution, score it with a proper rule''), not a side effect, but it is why the protocol adds metrics that neither objective optimises (Section 3) and reports where Pearson wins. A second circularity: choosing a judge or prompt by Pearson and then reporting Pearson; the protocol selects and reports with different data and different metrics (Section 6, selection arms).

\section*{3. Metrics, step by step}
For a held-out item $i$ with ratings $y_{ij}$, $j=1..k_i$, and forecast $q=p_i$. Every score is averaged over the ratings of the item, then over items in a held-out family, then over families. Lower is better unless stated.
\begin{itemize}
\item \textbf{RPS} (primary; normalised ranked probability score, range 0 to 1): $S(q,y)=\frac1{K-1}\sum_{k=0}^{K-2}\bigl(F_q(k)-\mathbf 1\{y\le k\}\bigr)^2$. Reference value: the uniform forecast. For ordered scales only.
\item \textbf{CE} (nats): $-\log q_\varepsilon(y)$, $q_\varepsilon=(1-\varepsilon)q+\varepsilon/K$, with $\varepsilon=6\times10^{-6}$ as in experiment 21, and the same table recomputed with $\varepsilon=10^{-2}$ (Sanmi's floor point). Reference: $\log K$ for uniform.
\item \textbf{TVD} (range 0 to 1): $\tfrac12\sum_r|q(r)-\hat h_i(r)|$ with $\hat h_i$ the item's empirical histogram. Descriptive when $k_i$ is small (R6). When an item has about 100 ratings (ChaosNLI, DICES) the histogram is close to $p^\ast_i$ and TVD and the expected scores below are meaningful.
\item \textbf{Top-1 agreement} (higher better): 1 if $\arg\max_r q(r)$ is among the human modal ratings of the item, else 0 (ties in the forecast share credit).
\item \textbf{Calibration of cumulative probabilities}: for each threshold $k$, sort held-out items by $F_q(k)$, cut into $B=5$ equal-count bins, and take the weighted mean $|\text{mean forecast}-\text{observed frequency of }y\le k|$; average over thresholds (expected calibration error, ECE). \textbf{Calibration slope:} least-squares fit $\hat t_i=a+b\,m_i$ on held-out items, where $\hat t_i$ is the human mean; ideal $(a,b)=(0,1)$. This measures what Pearson ignores.
\item \textbf{Pairwise preference accuracy} (higher better): over item pairs in the same family with human-mean gap $|\hat t_i-\hat t_j|\ge\delta$, the share where $\operatorname{sign}(m_i-m_j)=\operatorname{sign}(\hat t_i-\hat t_j)$, ties in $m$ counting $\tfrac12$; $\delta=1$ rating point (one point on a 0--5 scale; for $K=2$, any difference), fixed now. It reads means only, so it cannot separate shape; it is the metric for ranking uses.
\item \textbf{Secondary:} Pearson, Spearman, Lin's concordance correlation $\mathrm{CCC}=2s_{mt}/(s_m^2+s_t^2+(\bar m-\bar t)^2)$ (higher better).
\item \textbf{Target noise.} With one to three ratings per item the observed histogram is a poor stand-in for $p^\ast_i$. On datasets with about 100 ratings per item, train on a few ratings subsampled from each item but score against the full histogram (expected RPS and CE under $\hat h_i$). Both the sparse and the dense scorings are reported and labelled.
\end{itemize}

\section*{4. The ``who wins'' table (layout, to be filled by each run)}
One block per dataset, judge and training size $N$ (human ratings). Cell: difference with a 95\% interval and the winner. For lower-is-better metrics the difference is \emph{test metric of the comparator minus test metric of the CE-trained judge}; for higher-is-better metrics it is \emph{test metric of the CE-trained judge minus test metric of the comparator}. Positive always means the CE-trained judge is better. Absolute values of both baselines are printed beside it.

{\small\begin{longtable}{@{}>{\ra}p{3.0cm}>{\ra}p{2.5cm}>{\ra}p{3.0cm}>{\ra}p{2.5cm}>{\ra}p{2.5cm}@{}}
\toprule
Held-out metric & vs Pearson-trained judge & vs Pearson-trained judge with affine calibration & vs uniform forecast & vs training histogram\\
\midrule
\textbf{RPS} (primary) & $\Delta$ [lo, hi], winner & same & same & same\\
\textbf{CE}, floors $6\times10^{-6}$ and $10^{-2}$ & same & same & same & same\\
\textbf{TVD} & same & same & same & same\\
Top-1 agreement & same & same & same & same\\
Calibration ECE & same & same & same & same\\
Calibration slope (distance of $b$ from 1) & same & same & same & same\\
Pairwise preference accuracy & same & same & same & same\\
CCC, Spearman, Pearson (secondary) & same & same & same & same\\
\bottomrule
\end{longtable}}
Then a summary table: rows dataset $\times$ judge, columns $N^\ast$ under RPS, CE, TVD and under the neutral metrics, and a column ``where Pearson wins''.

\section*{5. Sanmi's points mapped to what the protocol does}
Sanmi's feedback of 11-16-2026 (saved in the repo, commit 10897ca). His five points, then the framing and prompts-vs-weights notes.
{\small\begin{longtable}{@{}>{\ra}p{3.3cm}>{\ra}p{5.0cm}>{\ra}p{7.2cm}@{}}
\toprule
Sanmi's point & Short form & What the protocol does\\
\midrule
1. Human data: the frontier mixture is null; both judges in the open-pool win lose to the constant forecast; single-judge selection favours correlation; add an RPS-selected arm & Report selection honestly & Selection arms: pick the best judge by Pearson, by CE and by RPS on the training ratings, and score the pick on held-out families. A comparison where both arms lose to uniform is labelled ``both worse than uniform'', never a win. CE reported at two floors.\\
2. Toy baselines: a Pearson-trained softmax pins only the mean, so it loses on RPS; the fair rival is score-then-calibrate; add affine-calibrated or MSE baselines; report RPS against the constant and the truth, not ratios of excess & Fair baselines and absolute scores & Arms include Pearson-trained with affine calibration (least-squares $a,b$ on training items, then exponential tilting of the judge's distribution to the calibrated mean) and an MSE-trained mean head with the same tilting. Every table gives absolute RPS next to uniform and (where known) the truth.\\
3. Small $n$: fitting many weights to a scale-free objective on 16 tasks overfits; the supportable claim is ``do not fit flexible models by maximising correlation'' & Smaller claim, regularisation & Training-size sweep from 4 to 256 ratings. Each arm gets its own ridge strength from one grid, chosen by inner leave-one-family-out cross-validation on that arm's own loss (no single shared ridge). Failures at small $N$ are reported (experiment 20: CE loses at 4 items).\\
4. Use: distribution accuracy is the target when the output is read as a rating or distribution; ranking uses need rank agreement; item signal vs rater severity & Metric follows use & Two use-cases, each with its own primary metric: distribution reading $\to$ RPS; ranking $\to$ pairwise accuracy and Spearman. Where raters have identifiers (DICES) fit an ordinal model with rater effects and score against the severity-adjusted expected rating as a second target; with two to three raters the mean is the target and the severity question is marked unidentified.\\
5. References: Lin's concordance, ICC consistency vs agreement, ordinal likelihood models with rater facets as proper-loss judges & Add agreement metrics and an ordinal arm & CCC is a secondary column; ICC (agreement form) between judge and human mean is reported when $\ge2$ raters per item; the ordinal-likelihood head is an extra proper-loss arm.\\
Framing & Narrower claim; drop ``don't need psychometrics'' & The claim under test is Sanmi's wording: correlation measures association with mean ratings and cannot certify agreement in level and scale, let alone spread; when output is read as a rating or distribution, select and report with a proper scoring rule such as RPS and keep correlation as a ranking diagnostic.\\
Prompts vs weights & Tune prompts and harness, coarse selection, run after the workshop deadline & A selection mode: choose among a small set of prompts or judges by Pearson, CE or RPS on $N$ ratings and report the selection regret on held-out families. Planned after the workshop deadline, as Sanmi suggested.\\
\bottomrule
\end{longtable}}

\section*{6. Arms, baselines, training sizes and intervals}
\textbf{Arms} (all fitted on the same training ratings, same features, same optimiser, same starting point): uniform forecast; training histogram with 0.5 pseudocount per bin; Pearson-trained; Pearson-trained with affine calibration; MSE-trained mean head with tilting; CE-trained; optionally RPS-trained and the ordinal-likelihood head. Judge weights are never changed when the judge is a frozen model: the fitted part is the small head (experiment 21 uses a 27-point grid of transformations and exhaustive search, so the optimum is guaranteed). Report the objective, search space and budget in every caption.

\textbf{Training size.} $N\in\{4,8,16,32,64,128,256\}$ human ratings, as in experiment 21, reported both as ratings and as items $\times$ ratings per item (one rating per item and two or three per item are separate rows). Experiment 20 found different thresholds for one and for three ratings per item.

\textbf{Splits and intervals.} Hold out whole task families (assignments, genres, prompt groups); at least 8 families. Per family, average the paired difference over 200 seeded training draws; then across the $F$ families the one-sided 97.5\% bound is mean $-\,t_{F-1,0.975}\cdot s\sqrt{1/F+1/(F-1)}$ (the overlap-adjusted form of experiment 21; $F=10$: 9 degrees of freedom), plus the $2^F$ sign-flip permutation check, family win counts, and variation across draws. Draws within a family are never counted as independent evidence. Judges are fixed in advance and $\alpha=0.05$ is split by their number.

\section*{7. Break-even rule: when is CE worth the cost of the ratings?}
For a dataset, judge and metric (RPS first), at each grid size $N$ let $d_N^{\rm P}$ and $d_N^{\rm U}$ be the held-out RPS of the Pearson-trained judge minus that of the CE-trained judge, and the same for the uniform forecast. $N$ \emph{passes} if the lower interval bound of both is above $0$. Then
\[
N^\ast=\min\{N\text{ in the grid}: N'\text{ passes for every grid size }N'\ge N\}
\]
(testing from the largest size down and stopping at the first failure, as in experiment 21; if no size passes, report ``$N^\ast>256$''). Report $N^\ast$ in ratings and in items.

\textbf{Cost.} Let $c$ be the cost of one human rating (money or minutes) and $B$ the budget, so $N_B=\lfloor B/c\rfloor$ ratings. \textbf{Rule:} if $N_B\ge N^\ast$, train with a proper loss; the break-even cost of the ratings is $C^\ast=c\,N^\ast$. If $N_B<N^\ast$, do not fit a flexible head by either loss: use the uniform forecast or the training histogram, or the Pearson-trained judge with affine calibration if it passes against uniform at $N_B$. Add the neutral-metric gate: CE must not lose to Pearson on top-1 agreement, calibration or pairwise accuracy beyond a tolerance fixed before the run (0.02 in each metric's units). Brando's rule, in his words (11-21-2026): ``If for the practical `cost' of getting human ratings/scores/judgements N is sufficiently large to use CE then we shouldn't be using Corr.''

\textbf{What is already measured} (experiment 20, softmax-linear head, synthetic, 3 ratings per item): CE beats the Pearson-trained judge from 16 items (48 ratings; +0.034, 96\% of seeds), ties at 8 items, loses at 4; with one rating per item it crosses at 32 items (32 ratings, 94\%). Real 75 artifacts: from 7 training families up, in the closed-weight 15, all 59 and historical 44 pools (+0.027 to +0.058); at 3 families it ties, and it loses in the open-weight 39 pool at 3 families ($-0.037$). These are not $N^\ast$ under this protocol (different splits and baselines).

\textbf{Experiment 21} (Codex, frozen before this note) already uses RPS first, uniform and Pearson baselines, 4--256 ratings and whole-assignment holdouts. This note does not change it. Top-1 agreement, calibration and pairwise accuracy can be computed from its saved forecasts at no cost; they would be labelled post hoc there.

\section*{8. Data checklist for public rating sets}
Before a set enters the protocol, record each item in a table next to the set's name:
\begin{enumerate}
\item \textbf{Individual ratings released} (not only the mean, majority or a filtered mean).
\item \textbf{Ratings per item} and whether it varies; enough to estimate a distribution (about $\ge5$) or to subsample from.
\item \textbf{Scale} $K$ and whether it is ordered (RPS needs order; for unordered labels use CE and TVD).
\item \textbf{Rater identity:} an anonymous id, the same raters across items, rater metadata; needed for rater-facet models.
\item \textbf{Item structure:} prompt, task or family id, to make family holdouts; whether pairs share a prompt.
\item \textbf{Licence:} commercial use, redistribution, attribution.
\item \textbf{Selection bias:} how items and raters were chosen (for example items picked for disagreement), and the rater quality filter.
\item \textbf{Same object as the judge sees:} the judge gets exactly what raters saw.
\item \textbf{Contamination:} possible exposure of public items to model pretraining (unknown unless stated).
\item \textbf{Sensitive content} and terms for model use.
\end{enumerate}

\textbf{Shortlist of public sets already checked} (read 11-21-2026 from each set's README or data card; DICES row counts computed from the released CSV files).
{\small\begin{longtable}{@{}>{\ra}p{2.0cm}>{\ra}p{3.9cm}>{\ra}p{2.6cm}>{\ra}p{2.6cm}>{\ra}p{3.2cm}@{}}
\toprule
Set & Releases & Scale, rater identity & Licence & Watch for\\
\midrule
ChaosNLI (\href{https://arxiv.org/abs/2010.03532}{paper}, \href{https://github.com/easonnie/ChaosNLI}{repo}) & 100 new labels for each of 4,645 examples: 3,113 from SNLI and MNLI, 1,532 from Abductive-NLI, as label counts per example (\texttt{label\_counter}), plus the earlier original labels; model logits for several models & $K=3$ (entailment, neutral, contradiction) for SNLI and MNLI; $K=2$ for Abductive-NLI. No rater id in the released format & CC BY-NC 4.0 (non-commercial), per the repo README & The SNLI and MNLI parts are only items whose majority label matched three of the five original labels (borderline agreement), so they lean toward disagreement; the Abductive-NLI part is the entire development set, so it is not selected. Workers were Mechanical Turk workers. With $K=2$ the mean fixes the whole distribution (W1), so Abductive-NLI can test level and scale only\\
DICES (\href{https://arxiv.org/abs/2306.11247}{paper}, \href{https://github.com/google-research-datasets/dices-dataset}{repo}) & DICES-350: 350 conversations, each rated by all 123 raters (43,050 rows); DICES-990: 990 conversations, 69 to 76 ratings each, 172 raters (72,103 rows). One row per (rater, conversation), with rater metadata, safety categories and a gold label & Yes / No / Unsure per category (unordered or weakly ordered); anonymous rater ids with demographics & CC BY 4.0 & Adversarial safety conversations with a contested notion of ground truth; unordered labels, so use CE and TVD, and RPS only with a stated ordering; the authors apply a rater-quality filter (104 of 123 raters in the 350 set) but release all\\
HelpSteer2 (\href{https://arxiv.org/abs/2406.08673}{paper}, \href{https://huggingface.co/datasets/nvidia/HelpSteer2}{data}) & 21,362 prompt-response samples with five attributes (helpfulness, correctness, coherence, complexity, verbosity). The main label is the mean of the three most agreeing annotators (3.41 annotators on average); a separate \texttt{disagreements} file (added 1 Oct 2024) gives the individual annotations as lists: 23,652 rows with 2 to 6 annotations each (1,162 with 2, 16,390 with 3, 1,314 with 4, 4,724 with 5, 62 with 6; counted from the file) & Integer 0 to 4, $K=5$, ordered; annotators are not identified & CC BY 4.0 & The main labels are filtered means, so use the \texttt{disagreements} lists; the number of annotations depends on difficulty (more annotators were added when the first three disagreed), so ratings per item are not independent of the item\\
\bottomrule
\end{longtable}}
Not yet surveyed: other public rating sets. The literature survey is a separate task, planned for Monday, and will fill this table with the same columns.

\section*{9. What this note does not do}
It runs nothing and cites no new results. It fixes definitions and decision rules on 11-21-2026; changes after any new data are looked at will be added as dated amendments, never as edits. Open points: the 23,652 rows of HelpSteer2's \texttt{disagreements} file do not equal the 21,362 main samples, so matching rows to main samples needs the prompt and response text; ChaosNLI's per-worker identity is not released, so rater-facet models are not possible there.
\end{document}
